% DRÖG — stutt lausn
Setjum $x = 0$:
\[
f(0) f(f(y)) = f(0).
\]
Ef $f(0) \neq 0$ fæst að $f(f(y)) = 1$ fyrir öll $y$.  
Það þýðir að $f$ myndi varpa öllum gildum á rauntölum þannig að samsett fall $f \circ f$ sé fast jafnt $1$.  
Setjum þetta í upphaflegu jöfnuna:
\[
f(x)(x + 1) = f(x^2) + x f(y).
\]
Við getum þá sett $y = 0$ og fáum
\[
f(x)(x + 1) = f(x^2) + x f(0).
\]
Frádráttur jafnanna gefur $x f(y) = x f(0)$ fyrir öll $x,y$. Fyrir $x \neq 0$ fæst $f(y) = f(0)$ fyrir öll $y$, sem þýðir að $f$ væri fastafall.  
Ef $f(x) \equiv c$, þá gefur upphaflega jafnan $c(x + c) = c x + x c$, sem einfaldast í $c^2 = c x$ fyrir öll $x$,  
sem er ómögulegt. Því hlýtur að gilda
\[
f(0) = 0.
\]

Setjum þá $y = 0$:
\[
f(x)\qty(x + f(f(0))) = f(x^2) + x f(0) \quad\Rightarrow\quad f(x)x = f(x^2).
\]
Þetta má setja inn í upphaflegu jöfnuna og fá:
\[
f(x) f(f(y)) = x f(y).
\]

Ef til er $a$ þannig að $f(a) \neq 0$, þá má setja $y = a$:
\[
f(x) f(f(a)) = x f(a).
\]
Ef $f(f(a)) = 0$, þá væri vinstri hlið jöfn núlli fyrir öll $x$, en hægri hliðin $x f(a)$ væri aðeins núll þegar $x = 0$,  
sem gefur mótsögn. Því hlýtur að gilda $f(f(a)) \neq 0$, og við getum deilt með því:
\[
f(x) = \frac{f(a)}{f(f(a))} \, x.
\]
Þannig fæst að $f(x) = c x$ fyrir einhvern fastann $c$.

Setjum $f(x) = c x$ inn í upphaflegu jöfnuna:
\[
c x \qty(x + c^2 y) = c x^2 + x c y,
\]
sem einfaldast í $c^3 = c$.

Því fæst
\[
c \in \{0, 1, -1\}.
\]

Almennar lausnir eru því
\[
f(x) = 0, \quad f(x) = x, \quad f(x) = -x.
\]

